\documentclass[10pt,a4paper]{article}

\usepackage[utf8]{inputenc}
\usepackage[T1]{fontenc}
\usepackage[spanish,es-tabla]{babel}

\usepackage{amsmath,amssymb,amsthm}
\usepackage{booktabs}
\usepackage{tikz}
\usetikzlibrary{automata,positioning,arrows}
\pagecolor{black!85!white}
\color{white}

\theoremstyle{definition}
\newtheorem{definicion}{Definición}[section]

\title{Hito 0 - Automatas finitos: enteros positivos}
\author{Jonathan Javier Pietracci}
\begin{document}

\section{Gramática de expresiones y el parser}

Para construir una calculadora necestamos reconocer expresiones como \texttt{2 + 3 * 4} o \texttt{(8-3) * 2}.
Esto requiere dos cosas: una \emph{gramática} que defina qué expresiones son válidas y un \emph{parser} que verifique si una cadena de tokes la satisface.

\subsection{Gramática ambigua(por qué no sirve?)}

La gramática ''obvia'' para expresiones es:

\begin{center}
\begin{tabular}{l}
    $E \to E + E$ \\
    $E \to E * E$ \\
    $E \to E$ \\
    $E to \text{número}$ \\
\end{tabular}
\end{center}

Esta gramática es \emph{ambigua}: existen cadenas con más de un árbol de derivación. Para \texttt{2 + 3 * 4} hay dos árboles posibles:

\begin{center}
\shorthandoff{<>}
\begin{tikzpicture}
    every node/.style={draw,circle,minimum size=0.7cm}
    \node{+}
     child { node {2} }
     child { node {$\times$} }
        child { node {3} }
        child { node {4} };
\end{tikzpicture}
\qquad\quad
\begin{tizkpicture}[level distance=1.2cm], sibling distance=1.5cm,
    every node/.style={draw,circle,minimum size=0.7cm}
\node {$times}
    child { node {+}
        child { node {2} }
        child { node {3} } }
    child { node {4} };
\end{tikzpicture}
\shorthandon{<>}
\end{center}

El de la izquierda evalúa $2 + (3 \times 4) = 14$ \\
El de la derecha evalúa $(2+3) \times 4 = 20$

Un compilador no puede elegir entre 2 árboles: necesitamos una gramática que produzca \emph{exactamente uno}.

\subsection{Gramática desambiguada (tres niveles)}

La solución es separar la precedencia en niveles de  no-terminales.
Cada nivel maneja una precedencia distinta, de menor (arriba) a mayor (abajo):

\begin{center}
\begin{tabular}{lll}
    \toprule
    Nivel & Reglas & Operadores \\
    \midrule
    $E$ (Expresión) & $E \to E + T \mid E - T \mid T$ & $+$ y $-$ \\
    $T$ (Término)   & $T \to T * F \mid T / F \mid F$ & $*$ y $/$ \\
    $F$ (Factor)    & $F \to (E) \mid \text{número}$  & paréntesis y números \\
    \bottomrule
\end{tabular}
\end{center}

La regla de oro: \textbf{cada nivel solo llama al nivel de abajo}. $E$ llama a $T$, $T$ llama a $F$, y $F$ llama a $E$ (solo para $\to$ multiplicación $\to$ suma.)

\subsection{Precedencia y asociatividad}

\begin{definicion}
La \textbf{asociatividad} determina el orden entre operadores de igual proecedencia. Si $E \to E + T$, el $E$ está a la izquierda, lo que fuerza asociatividad izquierda:
$(2 + 3) + 4$.
\end{definicion}

\begin{definicion}
    La \textbf{asociatividad} determina el orden entre operadores de igual precedencia. Si $E \to E + T$, el $E$ está
    a la izquierda, lo que fuerza asociatividad izquierda:
    $2 + 3 + 4$ se evalua como $(2 + 3) + 4$.
\end{definicion}


\subsection{Derivación de \texttt{2 + 3 * 4}}

Con la gramática desambiguada, la cadena \texttt{2 + 3 * 4} tiene una única derivación:


\end{document}
